\chapter{附录B：模型台账与敏感性}

\section{翻倍所需倍数}

\begin{exhibitbox}[Exhibit B1：现价隐含与翻倍门槛]
\small
\begin{tabularx}{\exhibitboxwidth}{L{2.4cm} R{1.6cm} R{1.6cm} R{1.7cm} R{1.7cm} X}
\toprule
标的 & EPS分母 & 当前隐含PE & 翻倍所需PE & 牛市目标 & 解释 \\
\midrule
雅化集团 & 2026E 2.63 & 6.38x & 12.77x & 42.08 & 盈利兑现后仍需倍数约翻倍 \\
恩捷股份 & 2027E 5.16 & 9.27x & 18.54x & 103.20 & 接近完整20x恢复估值 \\
中国船舶 & 2027E 3.13 & 10.55x & 21.10x & 68.9 & 质量高但重估幅度大 \\
江波龙 & 2026E 26.47 & 14.68x & 29.35x & 780 & 峰值EPS上给高倍数 \\
盛新锂能 & 2026E 2.45 & 11.29x & 22.57x & 59.0 & 商品周期与现金双重验证 \\
天华新能 & 2026E 5.70 & 9.80x & 19.61x & 102.4 & 牛市目标低于翻倍价 \\
\bottomrule
\end{tabularx}
\end{exhibitbox}

\section{正式目标权重与同源敏感性}

只有雅化集团进入正式权重：80\%经营驱动概率值、10\%市场现价、10\%外部目标，悲观/基准/牛市概率为30\%/50\%/20\%。外部42元与基准经营参数来自同一份东吴报告，不能被视作第二个独立证据；10\%权重只用于保留可审计的Street参照。若将外部权重降为零并按比例重分至经营驱动与现价，结果为23.47元；若采用70\%经营驱动、20\%现价、10\%外部目标，结果为24.57元；主模型为25.32元。恩捷103元及30.96/67.08/103.20元三档全部为同源零权敏感性。

\section{情景不是预测分布}

概率并非通过历史样本统计得到，而是将证据强弱、目标期限、现金流与经营桥压缩成透明判断。它们的作用是避免牛市目标“偷渡”为基准结论。若新证据出现，概率、EPS和倍数必须全部更新。

\section{模型边界}

报告没有为零权备选发布正式目标，也没有把券商目标反向修改成年末目标。天华新能即使盈利强，牛市区间仍不翻倍；江波龙即使数学上刚好翻倍，证据与现金流不支持正式目标；中国船舶即使质量最好，缺失点目标与高倍数要求仍然阻止升级；恩捷虽有点目标，但公司级驱动证据阻断估值资格。

\clearpage
